\section{System and Infrastructure}
\label{sec:infra}

Each interaction in MemT-1.0 depends on memory accumulated from earlier steps.
Training must combine these dependencies with batched GPU computation, while inference must maintain an independent evolving state for each trajectory.
Our infrastructure executes the model one decoder layer at a time, handles state recurrence within memory blocks, and manages model weights separately from trajectory-specific states.
We describe the resulting training schedule, memory operators, and online state organization below.

\subsection{Distributed Training}
\label{sec:distributed_training}

Training packs combine independent trajectories, each containing the steps defined in Section~\ref{sec:training}.
Attention stays within each step, while memory carries information between steps of the same trajectory and resets at trajectory boundaries.
Teacher forcing supplies all tokens in the pack, so the model processes the complete pack one decoder layer at a time.
Within each layer, attention and memory projections are batched across tokens.
The memory block alone advances state along each trajectory.
This organization retains large backbone matrix multiplications while preserving the order of memory writes and reads.

We use Megatron-Core~\cite{narayanan2021efficientlargescalelanguagemodel} and Bridge~\cite{nvidia2026megatronbridge} for distributed training, including expert parallelism for the backbone MoE layers.
Although backbone weights are frozen during memory training, gradients still pass through backbone computations and memory transitions, allowing later predictions to train earlier writes.
The remaining infrastructure challenge is therefore the cost of memory computation and its training activations, which we address below.

\subsection{Efficient Operators}
\label{sec:infra_operators}

Native memory introduces two related costs: ordered state updates and activations that grow with the number of memory experts.
Our execution separates state updates from reads of step snapshots, then groups complete steps into windows to limit the lifetime of expert projections.
This lets the memory branch control its working set while the backbone continues to process the full token pack layer by layer.

\textbf{Chunked updates and snapshot reads.}
The stable update has a rank-one structure.
Suppressing layer, step, and channel indices, let $\mathbf{k}_i$ denote the normalized key, and define $\mathbf{u}_i=\sqrt{\boldsymbol{\beta}_i}\odot\mathbf{k}_i$, $\mathbf{z}_i=\boldsymbol{\omega}_i\odot\mathbf{v}_i$, and $\mathbf{D}_i=\operatorname{Diag}(\exp(\boldsymbol{\lambda}_i))$.
The update in Section~\ref{sec:model_arc} can then be written as
\begin{equation*}
\mathbf{S}_i=\mathbf{D}_i\mathbf{S}_{i-1}
+\mathbf{u}_i^\top\bigl(\mathbf{z}_i-\mathbf{u}_i\mathbf{D}_i\mathbf{S}_{i-1}\bigr),
\qquad
\mathbf{S}_0=\mathbf{M}_t,\quad \mathbf{M}_{t+1}=\mathbf{S}_{L_t}.
\end{equation*}
Each token requires a vector--matrix multiplication and an outer product, with leading work of $O(d_kd_v)$ per channel.
We use chunked WY~\cite{NEURIPS2024_d13a3eae} to group these small dependent updates into triangular solves and matrix multiplications, carrying state between chunks and reconstructing intermediate quantities during backward.
The adaptation to MemT-1.0 is to pair this writer with a separate snapshot reader: all queries in step $t$ read the same incoming state $\mathbf{M}_t$, even as the writer processes several chunks to form $\mathbf{M}_{t+1}$.
The reads can therefore be batched independently, followed by RMSNorm and SoftCapSiLU output gating.
Gradients from later reads propagate through the writes that produced their snapshots.
Writer chunks determine computation granularity, and model steps determine when updated memory becomes visible.

\textbf{Complete-step windows for MoS.}
Dense MoS evaluates every expert, so materializing all expert projections at once requires tensors that grow with both pack length and expert count.
Whole-pack execution does exactly this: it projects every token in the pack before processing the resulting memory computations.
Windowed execution instead performs expert projections, state updates, and readout fusion for one group of complete steps at a time, passing a differentiable FP32 final state to the next window.
Both schedules use the same storage and utilization routers, shared expert addresses, and update kernel.
Storage probabilities modulate the erase gate, write gate, and log-decay. Utilization probabilities mix the gated expert readouts.

A window budget limits how many tokens are projected together, while keeping each step intact even if it exceeds the budget.
Recomputing internal intermediates during backward~\cite{korthikanti2023reducing} avoids retaining all windows' expert projections for the full training pass.
Smaller windows reduce live intermediates but require more launches and recomputation.
Outputs, routing probabilities, and boundary states still occupy memory, so the total remains dependent on pack length as well as window size.

\textbf{The memory--time tradeoff of window size.} We measure one complete MoS memory block on a single GPU using synthetic hidden-state inputs.
Table~\ref{tab:infra_window_tradeoff} fixes a 32K-token pack, BF16 parameters and projections, and 2048-token steps, comparing peak allocated memory and forward-plus-backward (F+B) time for two window sizes.
Arrows denote whole-pack execution $\rightarrow$ windowed execution; all values are medians.

\begin{table}[htbp]
    \centering
    \caption{Memory-time tradeoff for one MoS memory block on a single GPU at a 32K-token pack length.
    Arrows denote whole-pack to windowed execution. The values are medians.}
    \label{tab:infra_window_tradeoff}
    \small
    \setlength{\tabcolsep}{5pt}
    \renewcommand{\arraystretch}{1.15}
    \begin{tabular}{lrrrr}
        \toprule
        \makecell[l]{Window size\\(complete steps)} & \makecell{Peak memory\\(GiB)} & \makecell{Memory\\reduction} & \makecell{F+B\\(ms)} & \makecell{Time\\increase} \\
        \midrule
        2048 tokens (1 step) & $47.83 \rightarrow 6.03$ & 87.4\% & $462.7 \rightarrow 493.5$ & 6.6\% \\
        4096 tokens (2 steps) & $47.83 \rightarrow 9.49$ & 80.1\% & $463.0 \rightarrow 470.6$ & 1.6\% \\
        \bottomrule
    \end{tabular}
\end{table}

In this configuration, processing two steps at a time reduces peak memory by approximately 80\%, with a 1.6\% increase in median F+B time.
Shrinking the window to one step saves additional memory but raises the time increase to 6.6\%.
The window budget thus provides a direct resource tradeoff: smaller windows reduce live intermediates, while larger windows reduce execution overhead.
Appendix~\ref{appendix:infra_experiment_details} provides the block configuration, measurement protocol, and full results and variability from 8K to 32K.


\subsection{Stateful Inference and Serving}
\label{sec:infra_serving}

During online interaction, backbone and hyper memory weights remain fixed and are shared across trajectories, while each trajectory maintains its own dynamic memory states.
Within a step, generation reads the incoming memory snapshot, and attention over the current text uses that step's KV cache.
Once the interaction completes, storage incorporates the input and response into memory for the next interaction.
The serving runtime manages states by trajectory so that subsequent requests continue their respective histories.

For a fixed configuration, the number of memory-state elements per trajectory is
\begin{equation*}
N_{\mathrm{state}}=\sum_l E_l H_{kv,l}d_{k,l}d_{v,l},
\end{equation*}
where the sum covers memory layers and $E_l=1$ without MoS.
Growing history updates these elements while their number remains determined by configured memory capacity.
Total storage for multiple trajectories grows with trajectory count.
Active requests additionally require their current-step KV caches and computational workspaces.
Persistent-state storage can thus be estimated directly from the memory configuration, whereas runtime GPU memory also includes shared weights and active requests.
The effect of finite capacity on retained historical detail is examined in the capability evaluations and Discussion.

This organization separates history preparation from subsequent queries.
In Evaluation's two-GPU, single-request BEAM experiment, MemT-1.0-Nano's query latency after history preparation remains nearly constant as history grows.
At 64K tokens of history, MemT-1.0-Nano achieves an approximately \textbf{33$\times$ reduction} in logical per-trajectory state size compared with full context.
The storage metric measures the logical state size after history preparation and before querying, and excludes shared weights.
When subsequent queries reuse that state, the initial preparation cost is amortized over multiple interactions.
The full protocol and curves appear in Section~\ref{sec:evaluation} and Appendix~\ref{appendix:eval_storage_overhead}.
